\chapter[Posterior Manifold]{The Posterior Manifold}\label{ch:sec:the-posterior-manifold}

The Universal Optimality Defect measures distances in a Hilbert space
$V_I/\operatorname{Im}(C)$.  But that space captures the geometry of
incidence vectors, not of probability distributions themselves.
To connect the defect to actual encryption systems we need to embed the
simplex of posterior distributions into this geometry---to construct a
Riemannian manifold whose points are probability distributions and whose
distances coincide with the UOD.  This chapter builds that manifold,
starting from elementary Euclidean geometry.


The posterior manifold constructed in this chapter is the stage on which the entire geometric theory unfolds. Its Riemannian metric (Chapter~\ref{ch:sec:pullback-riemannian-metric}), the stability of the UOD (Chapter~\ref{ch:sec:stability-of-the-universal-optimality-defect}), and all subsequent bounds on information functionals (Chapters~\ref{ch:sec:entropy-bounds}--\ref{ch:sec:universal-security-bounds}) live on this manifold.

\section{Euclidean Preliminaries}\label{ch:sec:euclidean-preliminaries}
Let \(\mathbb{R}^{n}\) denote the Euclidean space endowed with its
standard inner product.
Define
$f(x)=\sum_{i=1}^{n} x_i.$
\begin{proposition}[Smoothness of the Sum]
\label{ch:prop:1-1}
The map \(f\) is smooth.
\end{proposition}
\begin{proof} The coordinate projections are linear. Finite sums of linear
maps are linear, and every linear map between finite-dimensional
Euclidean spaces is smooth. 
\end{proof}
\begin{proposition}[Differential of the Sum]
\label{ch:prop:1-2}
For every \(x\),
$Df_x(v)=\sum_i v_i.$
\end{proposition}
\begin{proof} Since \(f\) is linear, its derivative at any point equals
the map itself: \(Df_x(v)=f(v)=\sum_i v_i\).
\end{proof}
\begin{proposition}[Surjectivity]
\label{ch:prop:1-3}
The differential \(Df_x\) is surjective.
\end{proposition}
\begin{proof} For any \(t\in\mathbb R\) take \(v=(t,0,\dots,0)\).
Then \(Df_x(v)=t\), so the image of \(Df_x\) is all of \(\mathbb R\).
\end{proof}
\section{Interior Probability Simplex}\label{ch:sec:interior-probability-simplex}
Define
\[ \Delta_n^\circ=
\{p\in\mathbb R^{n}:p_i>0,\;\sum_i p_i=1\}. \]
\begin{theorem}[Smooth Manifold Structure]
\label{ch:thm:1-4}
The interior simplex is a smooth manifold of dimension \(n-1\).
\end{theorem}
\begin{proof} The affine hyperplane \[ \Delta_n^\circ=
\{p\in\mathbb R^{n}:p_i>0,\;\sum_i p_i=1\}. \] is an embedded manifold by
the Regular Level Set Theorem because \[ T_p\Delta_n^\circ=
\{v\in\mathbb R^{n}:\sum_i v_i=0\}. \] is surjective everywhere.
The interior simplex $\Delta_n^\circ$ is itself an open subset of this
hyperplane, so it inherits a smooth manifold structure. 
\end{proof}
\begin{theorem}[Tangent Space of the Simplex]\label{ch:thm:1-5}

For every \(p\in\Delta_n^\circ\),
\[ T_p\Delta_n^\circ=
\{v\in\mathbb R^{n}:\sum_i v_i=0\}. \]
\end{theorem}
\begin{proof} The simplex $\Delta_n^\circ$ is the regular level set
\(s^{-1}(1)\) of the sum map \(s(v)=\sum_i v_i\).  The tangent space
of a regular level set is the kernel of the differential, hence
\[
T_p\Delta_n^\circ=\ker(ds_p)=\{v\in\mathbb R^n:\textstyle\sum_i v_i=0\}.
\]
\end{proof}

\begin{proposition}[Gauge Geometry of the Logarithmic Space]\label{ch:prop:2-1}
Define

\[
x\sim y \iff \exists\lambda\in\mathbb R,\; y=x+\lambda\mathbf1.
\]

This is an equivalence relation.  The quotient space

\[
\mathcal G = \mathbb R^{k}/\operatorname{span}\{\mathbf1\}
\]

inherits the canonical vector-space structure.
\end{proposition}

\begin{proposition}\label{ch:prop:2-1-dim}
\(\dim(\mathcal G)=k-1.\)
\end{proposition}
\begin{proof} Apply the dimension formula for quotient spaces.
\end{proof}


\section{Canonical Projection}\label{ch:sec:canonical-projection}

Define the quotient map

\[
Q:\mathbb R^{k}\to\mathcal G,\qquad Q(x)=[x].
\]

\begin{proposition}
\label{ch:prop:2-2}
The map \(Q\) is linear.
\end{proposition}

\begin{proposition}
\label{ch:prop:2-3}
\[ \ker(Q)=\operatorname{span}\{\mathbf1\}. \]
\end{proposition}
\begin{corollary}
\label{ch:cor:2-4}
\[ \operatorname{rank}(Q)=k-1. \]
\end{corollary}

\section{Hyperplane Representation}\label{ch:sec:hyperplane-representation}

Define

\[
H=\{x:\sum_i x_i=0\}.
\]

\begin{proposition}
\label{ch:prop:2-5}
The restriction \(Q|_H:H\to\mathcal G\) is a linear isomorphism.
\end{proposition}

Thus the quotient and hyperplane realizations are equivalent.

\paragraph{Logarithmic Coordinates.}
Define the logarithmic embedding

\[
L(P)=(\log P_1,\ldots,\log P_k).
\]

The map is smooth.  Its differential is

\[
DL_P(v)=\left(\frac{v_i}{P_i}\right)_i.
\]

\section{Weighted Hyperplane}\label{ch:sec:weighted-hyperplane}

Define

\[
W_P = \left\{ w:\sum_iP_iw_i=0 \right\}.
\]

\begin{proposition}[Weighted Hyperplane]
\label{ch:prop:3-1}
\(W_P\) is a vector subspace of dimension \(k-1\).
\end{proposition}

\paragraph{Interpretation.}
The posterior simplex $\Delta_k^\circ$ is the space of all possible posterior distributions.  Endowing it with a Riemannian metric turns probabilistic statements into geometric ones: distances become indistinguishability, geodesics become optimal interpolation, and curvature becomes sensitivity.  This geometric recasting is what makes the universal bounds of Part~IV possible.

\begin{quote}
\textbf{Mathematical intuition.}
The weighted hyperplane $W_P=\{w:\sum_i P_i w_i=0\}$ removes the degree of freedom that corresponds to multiplying all probabilities by the same factor.
This is the logarithmic analogue of normalising a probability distribution to sum to~1, but done in the tangent space rather than on the simplex itself.
\end{quote}

\begin{theorem}[Image of the Differential]
\label{ch:thm:3-2}
\[ DL_P(T_P\Delta)=W_P. \]
\end{theorem}

\begin{proof}
If \(v\in T_P\Delta\) then \(\sum_i v_i=0\).  Hence
\(\sum_i P_i(v_i/P_i)=\sum_i v_i=0\), so \(DL_P(v)\in W_P\).
Conversely, given \(w\in W_P\), define \(v_i=P_iw_i\).  Then
\(\sum_i v_i=\sum_i P_iw_i=0\), so \(v\in T_P\Delta\) and
\(DL_P(v)=w\). \(\)
\end{proof}

\begin{corollary}[Linear Isomorphism]
\label{ch:cor:3-3}
\(DL_P:T_P\Delta\to W_P\) is a linear isomorphism.
\end{corollary}

\paragraph{Running example: differential privacy.}
In differential privacy (Section~\ref{ch:sec:differential-privacy-and-database-queries}), the posterior manifold is the set of all output distributions induced by neighbouring databases.  The logarithmic map $L(P)=\log P$ sends the simplex to a hyperplane in $\mathbb R^n$, and the differential $DL_P$ gives the linearised effect of a database change.  For the Laplace mechanism with $n$ possible outputs, the weighted hyperplane $W_P=\{x\in\mathbb R^n:\sum_i\pi_i x_i=0\}$ has dimension $n-1$, and $\dim(\operatorname{Im}(DL_P))=k-1$ by Corollary~\ref{ch:cor:3-3}---every privacy variation is captured by a unique logarithmic direction.

This weighted hyperplane is the natural codomain of the logarithmic
differential and provides the key structural ingredient for the
subsequent geometry.
